\documentclass[12pt, titlepage]{article}

\usepackage{fullpage}
\usepackage[round]{natbib}
\usepackage{multirow}
\usepackage{booktabs}
\usepackage{tabularx}
\usepackage{graphicx}
\usepackage{float}
\usepackage{hyperref}
\hypersetup{
    colorlinks,
    citecolor=blue,
    filecolor=black,
    linkcolor=red,
    urlcolor=blue
}

\input{../../Comments.text}
\input{../../Common.text}

\newcounter{acnum}
\newcommand{\actheacnum}{AC\theacnum}
\newcommand{\acref}[1]{AC\ref{#1}}

\newcounter{ucnum}
\newcommand{\uctheucnum}{UC\theucnum}
\newcommand{\uref}[1]{UC\ref{#1}}

\newcounter{mnum}
\newcommand{\mthemnum}{M\themnum}
\newcommand{\mref}[1]{M\ref{#1}}

\begin{document}

\title{Module Guide for \progname{}} 
\author{\authname}
\date{\today}

\maketitle

\pagenumbering{roman}

\section{Revision History}

\begin{tabularx}{\textwidth}{p{3cm}p{2cm}X}
\toprule {\bf Date} & {\bf Version} & {\bf Notes}\\
\midrule
Date 1 & 1.0 & Notes\\
Date 2 & 1.1 & Notes\\
\bottomrule
\end{tabularx}

\newpage

\section{Reference Material}

This section records information for easy reference. 

\subsection{Abbreviations and Acronyms}

\renewcommand{\arraystretch}{1.2}
\begin{tabular}{l l} 
  \toprule		
  \textbf{symbol} & \textbf{description}\\
  \midrule 
  AC & Anticipated Change\\
  DAG & Directed Acyclic Graph \\
  M & Module \\
  MG & Module Guide \\
  OS & Operating System \\
  R & Requirement\\
  SC & Scientific Computing \\
  SRS & Software Requirements Specification\\
  \progname & Explanation of program name\\
  UC & Unlikely Change \\
  \wss{etc.} & \wss{...}\\
  \bottomrule
\end{tabular}\\

\newpage

\tableofcontents

\listoftables

\listoffigures

\newpage

\pagenumbering{arabic}

\section{Introduction}

Decomposing a system into modules is a commonly accepted approach to developing
software.  A module is a work assignment for a programmer or programming
team~\citep{ParnasEtAl1984}.  We advocate a decomposition
based on the principle of information hiding~\citep{Parnas1972a}.  This
principle supports design for change, because the ``secrets'' that each module
hides represent likely future changes.  Design for change is valuable in SC,
where modifications are frequent, especially during initial development as the
solution space is explored.  

Our design follows the rules layed out by \citet{ParnasEtAl1984}, as follows:
\begin{itemize}
\item System details that are likely to change independently should be the
  secrets of separate modules.
\item Each data structure is implemented in only one module.
\item Any other program that requires information stored in a module's data
  structures must obtain it by calling access programs belonging to that module.
\end{itemize}

After completing the first stage of the design, the Software Requirements
Specification (SRS), the Module Guide (MG) is developed~\citep{ParnasEtAl1984}. The MG
specifies the modular structure of the system and is intended to allow both
designers and maintainers to easily identify the parts of the software.  The
potential readers of this document are as follows:

\begin{itemize}
\item New project members: This document can be a guide for a new project member
  to easily understand the overall structure and quickly find the
  relevant modules they are searching for.
\item Maintainers: The hierarchical structure of the module guide improves the
  maintainers' understanding when they need to make changes to the system. It is
  important for a maintainer to update the relevant sections of the document
  after changes have been made.
\item Designers: Once the module guide has been written, it can be used to
  check for consistency, feasibility, and flexibility. Designers can verify the
  system in various ways, such as consistency among modules, feasibility of the
  decomposition, and flexibility of the design.
\end{itemize}

The rest of the document is organized as follows. Section
\ref{SecChange} lists the anticipated and unlikely changes of the software
requirements. Section \ref{SecMH} summarizes the module decomposition that
was constructed according to the likely changes. Section \ref{SecMD} gives a
detailed description of the modules. Section \ref{SecConnection} specifies the
connections between the software requirements and the modules, and
Section \ref{SecTM} within it includes two traceability matrices. One checks
the completeness of the design against the requirements provided in the SRS. The
other shows the relation between anticipated changes and the modules. Section
\ref{SecUse} describes the use relation between modules.

\section{Anticipated and Unlikely Changes} \label{SecChange}

This section lists possible changes to the system. According to the likeliness
of the change, the possible changes are classified into two
categories. Anticipated changes are listed in Section \ref{SecAchange}, and
unlikely changes are listed in Section \ref{SecUchange}. The information for the sections was based 
upon the initial findings stated in the \href{URL}{SRS}.

\subsection{Anticipated Changes} \label{SecAchange}

\begin{description}
  \item[AC1: Data collection sources.] The sources from which fitness data can be tracked will be changed and updated fitness data providers are added or update their interfaces.
  \item[AC2: Educational content.] The content of the financial lessons will change as technical decisions are made.
  \item[AC3: Virtual reward system.] The rules for calculating reward values and amounts will be adjusted based on application testing to ensure the best user experience.
  \item[AC4: User interface look and feel.] The user interface will change over time as development continues, in order to provide the best user experience.
  \item[AC5: Market simulator.] The generation of assets and securities, and their rates of growth and decline, will require calibration to provide a realistic experience while balancing user time constraints and user experience.
  \item[AC6: Notifications.] The kinds of notifications and their default settings are likely to change based on user feedback.
  \item[AC7: Hardware.] The specific hardware on which the software runs, including the phone's sensors and the wearable devices that record user activity, may change as new devices are supported.
  \item[AC9: Sign-in methods.] The sign-in methods supported are likely to expand, for example through the addition of sign-in with a third-party identity provider such as Google.
  \item[AC10: Fund pricing.] The method by which ETF and fund prices are determined may change, for example by deriving them from the simulated assets each fund holds instead of simulating them independently.
\end{description}

\subsection{Unlikely Changes} \label{SecUchange}

\begin{description}
  \item[UC1: Core functionality.] The following core functions of the application will not change:
    \begin{itemize}
      \item \textbf{Investing:} Users can purchase and sell securities and assets.
      \item \textbf{Portfolio:} Users can hold securities, and the portfolio's value can grow or decline.
      \item \textbf{Fitness goals:} Users can set and work toward fitness goals to earn virtual currency.
      \item \textbf{Market simulation:} The market is a simulated environment that is separate from the real world.
    \end{itemize}
  \item[UC2: Virtual currency.] The virtual currency used in the application has no monetary value.
  \item[UC3: Social features.] The following social features will not change:
    \begin{itemize}
      \item \textbf{Friends:} Users can add friends and communicate with them.
      \item \textbf{Leaderboards:} Users can join groups and compete against each other on multiple statistics, such as portfolio performance and currency acquired.
    \end{itemize}
  \item[UC4: Privacy.] User data will be protected in accordance with all applicable legal requirements.
  \item[UC5: Platform.] The application will be developed for the Android operating system. Additional platforms will not be added due to time constraints.
  \item[UC6: Technology stack.] The tools used to build the application will remain consistent throughout development.
  \item[UC7: Web API.] The API will be developed as a web API.
\end{description}

\section{Module Hierarchy} \label{SecMH}

This section provides an overview of the module design. Modules are summarized
in a hierarchy decomposed by secrets in Table \ref{TblMH}. The modules listed
below, which are leaves in the hierarchy tree, are the modules that will
actually be implemented.

The mobile application and the backend service are treated as a single module
hierarchy rather than two separate subsystems. The boundary between them is the
API Client Module (\mref{mAPI}): mobile modules reach backend behaviour only
through it, so the transport between the two can change without affecting
either side.

\begin{description}
\item [\refstepcounter{mnum} \mthemnum \label{mFI}:] Fitness Integration Module
\item [\refstepcounter{mnum} \mthemnum \label{mPer}:] Persistence Module
\item [\refstepcounter{mnum} \mthemnum \label{mUser}:] User and Onboarding Module
\item [\refstepcounter{mnum} \mthemnum \label{mAct}:] Activity Module
\item [\refstepcounter{mnum} \mthemnum \label{mGoal}:] Goals Module
\item [\refstepcounter{mnum} \mthemnum \label{mCur}:] Currency Ledger Module
\item [\refstepcounter{mnum} \mthemnum \label{mAsset}:] Asset Module
\item [\refstepcounter{mnum} \mthemnum \label{mPurch}:] Transaction Order Module
\item [\refstepcounter{mnum} \mthemnum \label{mPort}:] Portfolio Module
\item [\refstepcounter{mnum} \mthemnum \label{mLead}:] Leaderboard and Friends Module
\item [\refstepcounter{mnum} \mthemnum \label{mProg}:] Progress Reporting Module
\item [\refstepcounter{mnum} \mthemnum \label{mUI}:] UI Screens Module
\item [\refstepcounter{mnum} \mthemnum \label{mNoti}:] Notification Module
\item [\refstepcounter{mnum} \mthemnum \label{mMkt}:] Market Simulation Module
\item [\refstepcounter{mnum} \mthemnum \label{mAPI}:] API Client Module
\item [\refstepcounter{mnum} \mthemnum \label{mAuth}:] Authentication and Security Module
\end{description}


\begin{table}[H]
\centering
\begin{tabular}{p{0.3\textwidth} p{0.6\textwidth}}
\toprule
\textbf{Level 1} & \textbf{Level 2}\\
\midrule

Hardware-Hiding Module & Fitness Integration Module\\
\midrule

\multirow{14}{0.3\textwidth}{Behaviour-Hiding Module} & Persistence Module\\
& User and Onboarding Module\\
& Activity Module\\
& Goals Module\\
& Currency Ledger Module\\
& Asset Module\\
& Transaction Order Module\\
& Portfolio Module\\
& Leaderboard and Friends Module\\
& Progress Reporting Module\\
& UI Screens Module\\
& Notification Module\\
& Market Simulation Module\\
& API Client Module\\
\midrule

Software Decision Module & Authentication and Security Module\\
\bottomrule

\end{tabular}
\caption{Module Hierarchy}
\label{TblMH}
\end{table}

\section{Module Decomposition} \label{SecMD}

Modules are decomposed according to the principle of ``information hiding''
proposed by \citet{ParnasEtAl1984}. The \emph{Secrets} field in a module
decomposition is a brief statement of the design decision hidden by the
module. The \emph{Services} field specifies \emph{what} the module will do
without documenting \emph{how} to do it. For each module, a suggestion for the
implementing software is given under the \emph{Implemented By} title. If the
entry is \emph{OS}, this means that the module is provided by the operating
system or by standard programming language libraries.  \emph{\progname{}} means the
module will be implemented by the \progname{} software.

Only the leaf modules in the hierarchy have to be implemented. If a dash
(\emph{--}) is shown, this means that the module is not a leaf and will not have
to be implemented.

\subsection{Hardware-Hiding Modules}

\begin{description}
\item[Secrets:] How the system interacts with physical devices: the phone's
  sensors and the wearables that record the user's activity, and the platform
  services that expose their readings.
\item[Services:] Serves as a virtual layer over the devices that physically
  measure the user's activity, so the rest of the system can read activity data
  without knowing which device recorded it or how that device is accessed. Only
  modules that directly manipulate or read from hardware belong here; modules
  that merely use storage or the network, where the system does not choose or
  control the physical device, are behaviour-hiding.
\item[Implemented By:] --
\end{description}

\subsubsection{Fitness Integration Module (\mref{mFI})}

\begin{description}
\item[Secrets:] How each fitness data provider (Health Connect, Strava, Garmin)
  is accessed: its SDK or web API, its permission model, its record formats,
  how manually entered records are recognised and excluded, and how step counts
  recorded by several sources are combined without double counting.
\item[Services:] Reports whether a provider is available on the device,
  requests read permission for the data the system needs, and returns workouts
  and daily step totals since a given date in a common, provider-independent
  format. Each record carries a stable key so that the same activity is never
  reported as new twice.
\item[Implemented By:] \progname{} (mobile), using the provider SDKs and APIs
\item[Type of Module:] Abstract Data Type (one instance per provider, all
  sharing a single interface)
\end{description}

\subsection{Behaviour-Hiding Module}

\begin{description}
\item[Secrets:] The contents of the required behaviours, and how the data and
  services those behaviours depend on are represented within \progname{}.
\item[Services:] Includes programs that provide externally visible behaviour of
  the system as specified in the software requirements specification (SRS)
  documents, together with the modules that are specific to \progname{} and
  support that behaviour: storing the system's entities, simulating the market
  and communicating between the mobile app and the backend. Storage is
  behaviour-hiding rather than hardware-hiding because the system does not
  choose or manipulate where its data is physically stored; it only decides
  what is stored and how it is organised. The programs in this module will need
  to change if there are changes in the SRS.
\item[Implemented By:] --
\end{description}

\subsubsection{Persistence Module (\mref{mPer})}

\begin{description}
\item[Secrets:] The database management system, the database schema, and how
  entities are mapped to tables, queried and migrated between versions.
\item[Services:] Stores and retrieves the system's persistent entities (users,
  activities, goals, ledger entries, assets, prices, holdings, transactions,
  leaderboard groups and notification preferences) on behalf of the backend
  modules.
\item[Implemented By:] \progname{} (backend), using Spring Data JPA, Liquibase
  and MySQL
\item[Type of Module:] Abstract Object
\end{description}

\subsubsection{User and Onboarding Module (\mref{mUser})}

\begin{description}
\item[Secrets:] How user accounts and onboarding questionnaire responses are
  represented and validated.
\item[Services:] Registers new users with a unique email address, looks up and
  updates user profiles, and records the user's onboarding answers about their
  fitness level and financial background.
\item[Implemented By:] \progname{} (backend)
\item[Type of Module:] Abstract Object
\end{description}

\subsubsection{Activity Module (\mref{mAct})}

\begin{description}
\item[Secrets:] How synced activity is represented, stored and deduplicated,
  and which records are eligible for a reward.
\item[Services:] Accepts workouts and daily step totals synced from the mobile
  app, ignores records that have already been received, reports a user's
  recorded activity over a given period, and tells the Currency Ledger Module
  when new eligible activity should be rewarded.
\item[Implemented By:] \progname{} (backend)
\item[Type of Module:] Abstract Object
\end{description}

\subsubsection{Goals Module (\mref{mGoal})}

\begin{description}
\item[Secrets:] How fitness goals are represented and how progress towards a
  goal is measured from recorded activity.
\item[Services:] Lets users create, update and cancel fitness goals, measures
  progress towards each active goal from the activity reported by the Activity
  Module, marks a goal complete when its target is reached, and tells the
  Currency Ledger Module when a completed goal should be rewarded.
\item[Implemented By:] \progname{} (backend)
\item[Type of Module:] Abstract Object
\end{description}

\subsubsection{Currency Ledger Module (\mref{mCur})}

\begin{description}
\item[Secrets:] The reward rule that converts activity into virtual currency
  (rates, bonuses and caps), and how the ledger of credits and debits is
  represented.
\item[Services:] Credits currency for newly rewarded activity and goals, debits
  currency for purchases, credits the proceeds of sales, and reports each
  user's current balance and history of entries. A balance can never become
  negative, and currency can only be created by fitness activity or investment
  returns.
\item[Implemented By:] \progname{} (backend)
\item[Type of Module:] Abstract Object
\end{description}

\subsubsection{Asset Module (\mref{mAsset})}

\begin{description}
\item[Secrets:] How the simulated assets and asset types (stocks, bonds, ETFs
  and funds) are defined and represented, and how many units of each asset are
  available in the simulated market.
\item[Services:] Lists the assets that can be traded, describes each asset
  (symbol, name and asset type), and reports and updates the quantity of each
  asset available. A new asset type can be added through new asset definitions
  alone, without changes to the trading logic.
\item[Implemented By:] \progname{} (backend)
\item[Type of Module:] Abstract Object
\end{description}

\subsubsection{Transaction Order Module (\mref{mPurch})}

\begin{description}
\item[Secrets:] How buy and sell orders are validated and executed, and how the
  record of executed transactions is represented.
\item[Services:] Executes buy and sell orders at the asset's current simulated
  price, rejects orders the user cannot afford or holdings they do not own,
  ensures a repeated order request is executed only once, debits or credits the
  user's balance through the Currency Ledger Module, updates the user's
  holdings through the Portfolio Module, and records a transaction for every
  executed trade. An order either completes in full or has no effect.
\item[Implemented By:] \progname{} (backend)
\item[Type of Module:] Abstract Object
\end{description}

\subsubsection{Portfolio Module (\mref{mPort})}

\begin{description}
\item[Secrets:] How a user's holdings and their cost are represented, and how
  portfolio value and gains or losses are calculated.
\item[Services:] Maintains one portfolio per user, reports the quantity of each
  asset the user holds, and reports the portfolio's total value and its gains
  or losses at current simulated prices. Holdings change only when the
  Transaction Order Module executes a trade.
\item[Implemented By:] \progname{} (backend)
\item[Type of Module:] Abstract Object
\end{description}

\subsubsection{Leaderboard and Friends Module (\mref{mLead})}

\begin{description}
\item[Secrets:] How friend lists and leaderboard groups are stored, and how
  members of a group are ranked.
\item[Services:] Lets users manage a list of friends, create private
  leaderboard groups and invite friends to them, join and leave public or
  private groups, and view each group's members ranked by progress towards
  their goals. Leaderboard standings never award currency.
\item[Implemented By:] \progname{} (backend)
\item[Type of Module:] Abstract Object
\end{description}

\subsubsection{Progress Reporting Module (\mref{mProg})}

\begin{description}
\item[Secrets:] Which figures are combined into a progress summary and how they
  are aggregated over a reporting period.
\item[Services:] Produces summaries of a user's fitness activity, goal
  completion, currency earned, portfolio performance and leaderboard standings
  by querying the modules that own that data.
\item[Implemented By:] \progname{} (backend)
\item[Type of Module:] Library
\end{description}

\subsubsection{UI Screens Module (\mref{mUI})}

\begin{description}
\item[Secrets:] The layout and visual style of each screen (colours,
  typography and spacing), the navigation between screens, and
  how loading, empty and error states are presented to the user.
\item[Services:] Presents the sign-in and onboarding flows, the dashboard, the
  activity screen, the exchange and asset screens, the order screen and the
  profile screen, and turns user input into calls on the Fitness Integration
  and API Client modules.
\item[Implemented By:] \progname{} (mobile), using React Native and Expo Router
\item[Type of Module:] Library
\end{description}

\subsubsection{Notification Module (\mref{mNoti})}

\begin{description}
\item[Secrets:] Which events produce a notification, how each message is
  worded, how the user's notification preferences, quiet hours and daily limit
  are enforced, and how notifications are delivered to the device.
\item[Services:] Sends workout reminders, reward and achievement notices and
  simulated market notices to the user. It lets users turn each kind of
  notification on or off and set quiet hours. It sends nothing during quiet
  hours or for a kind the user has turned off, and by default sends no more
  than \textsf{MaxDailyNotifications} goal-related notifications per day. If
  notifications are disabled or delivery fails, the failure is handled without
  affecting the rest of the system.
\item[Implemented By:] \progname{} (backend and mobile), using the Android
  notification service
\item[Type of Module:] Abstract Object
\end{description}

\subsubsection{Market Simulation Module (\mref{mMkt})}

\begin{description}
\item[Secrets:] The algorithm used to simulate price movement over time, how
  it is calibrated for each asset type, how market-wide downturns are
  generated, and how dividends are scheduled, sized and paid.
\item[Services:] Provides the current price of each asset defined by the Asset
  Module and its price history, and updates prices on a schedule. Reports a
  trend for each asset type and the market events that have occurred, such as
  downturns and dividends. Announces each asset's dividends and pays them to
  the users holding that asset through the Currency Ledger Module. Generates a
  past price history for a newly listed asset.
\item[Implemented By:] \progname{} (backend)
\item[Type of Module:] Abstract Object
\end{description}

\subsubsection{API Client Module (\mref{mAPI})}

\begin{description}
\item[Secrets:] The backend's address, the transport protocol and message
  format, how credentials and the device identifier are stored and attached
  to requests, how network failures are retried or reported, how received
  data is kept for use when the backend cannot be reached, and how a repeated
  request is prevented from taking effect twice.
\item[Services:] Gives the mobile modules one typed function per backend
  operation (for example syncing activity, reading the balance, listing assets
  and placing an order), so they never build requests or parse responses
  themselves. Keeps the user signed in by renewing access when it expires,
  returns the last data received, marked as possibly out of date, when the
  backend cannot be reached, and reports failures as typed errors that the
  screens can explain to the user.
\item[Implemented By:] \progname{} (mobile)
\item[Type of Module:] Abstract Object
\end{description}

\subsection{Software Decision Module}

\begin{description}
\item[Secrets:] General-purpose design decisions based on mathematical
  theorems, established algorithms or programming considerations. The secrets
  of these modules are \emph{not} described in the SRS and do not depend on
  anything specific to \progname{}.
\item[Services:] Includes abstract data structures and algorithms that do not
  provide direct interaction with the user and that could be reused, largely
  unchanged, in other projects.
  % Changes in these modules are more likely to be motivated by a desire to
  % improve performance than by externally imposed changes.
\item[Implemented By:] --
\end{description}

\subsubsection{Authentication and Security Module (\mref{mAuth})}

\begin{description}
\item[Secrets:] The authentication mechanism, how credentials are hashed and
  verified, how sessions or tokens are issued, which requests require an
  authenticated user, how multi-factor authentication codes and password reset
  links are generated and sent by email, how request rates are limited, and how
  security events are logged.
\item[Services:] Verifies a user's credentials, enforces the password policy
  and the sign-in lockout, offers optional multi-factor authentication and
  password reset through the user's email, identifies the user making each
  request to the backend, and rejects requests that are unauthenticated, that
  exceed the request rate limit, or that try to access another user's data.
  Ends sessions when the user signs out, after a period without use, and when
  the account is deleted. Records failed sign-ins, rejected tokens, refused
  access and rate limit violations in the security log.
\item[Implemented By:] \progname{} (backend), using Spring Security
\item[Type of Module:] Abstract Object
\end{description}

\section{Connection Between Requirements and Design} \label{SecConnection}

The design of the system is intended to satisfy the requirements developed in
the SRS. In this stage, the system is decomposed into modules. The connection
between requirements and modules is listed in Table~\ref{TblRT}.

\wss{The intention of this section is to document decisions that are made
  ``between'' the requirements and the design.  To satisfy some requirements,
  design decisions need to be made.  Rather than make these decisions implicit,
  they are explicitly recorded here.  For instance, if a program has security
  requirements, a specific design decision may be made to satisfy those
  requirements with a password.}

\subsection{Traceability Matrix} \label{SecTM}

This section shows two traceability matrices: between the modules and the
requirements and between the modules and the anticipated changes.

% the table should use mref, the requirements should be named, use something
% like fref
\begin{table}[H]
\centering
\begin{tabular}{p{0.2\textwidth} p{0.6\textwidth}}
\toprule
\textbf{Req.} & \textbf{Modules}\\
\midrule
R1 & \mref{mHH}, \mref{mInput}, \mref{mParams}, \mref{mControl}\\
R2 & \mref{mInput}, \mref{mParams}\\
R3 & \mref{mVerify}\\
R4 & \mref{mOutput}, \mref{mControl}\\
R5 & \mref{mOutput}, \mref{mODEs}, \mref{mControl}, \mref{mSeqDS}, \mref{mSolver}, \mref{mPlot}\\
R6 & \mref{mOutput}, \mref{mODEs}, \mref{mControl}, \mref{mSeqDS}, \mref{mSolver}, \mref{mPlot}\\
R7 & \mref{mOutput}, \mref{mEnergy}, \mref{mControl}, \mref{mSeqDS}, \mref{mPlot}\\
R8 & \mref{mOutput}, \mref{mEnergy}, \mref{mControl}, \mref{mSeqDS}, \mref{mPlot}\\
R9 & \mref{mVerifyOut}\\
R10 & \mref{mOutput}, \mref{mODEs}, \mref{mControl}\\
R11 & \mref{mOutput}, \mref{mODEs}, \mref{mEnergy}, \mref{mControl}\\
\bottomrule
\end{tabular}
\caption{Trace Between Requirements and Modules}
\label{TblRT}
\end{table}

\begin{table}[H]
\centering
\begin{tabular}{p{0.2\textwidth} p{0.6\textwidth}}
\toprule
\textbf{AC} & \textbf{Modules}\\
\midrule
\acref{acHardware} & \mref{mHH}\\
\acref{acInput} & \mref{mInput}\\
\acref{acParams} & \mref{mParams}\\
\acref{acVerify} & \mref{mVerify}\\
\acref{acOutput} & \mref{mOutput}\\
\acref{acVerifyOut} & \mref{mVerifyOut}\\
\acref{acODEs} & \mref{mODEs}\\
\acref{acEnergy} & \mref{mEnergy}\\
\acref{acControl} & \mref{mControl}\\
\acref{acSeqDS} & \mref{mSeqDS}\\
\acref{acSolver} & \mref{mSolver}\\
\acref{acPlot} & \mref{mPlot}\\
\bottomrule
\end{tabular}
\caption{Trace Between Anticipated Changes and Modules}
\label{TblACT}
\end{table}

\section{Use Hierarchy Between Modules} \label{SecUse}

In this section, the uses hierarchy between modules is
provided. \citet{Parnas1978} said of two programs A and B that A {\em uses} B if
correct execution of B may be necessary for A to complete the task described in
its specification. That is, A {\em uses} B if there exist situations in which
the correct functioning of A depends upon the availability of a correct
implementation of B.  Figure \ref{FigUH} illustrates the use relation between
the modules. It can be seen that the graph is a directed acyclic graph
(DAG). Each level of the hierarchy offers a testable and usable subset of the
system, and modules in the higher level of the hierarchy are essentially simpler
because they use modules from the lower levels.

\wss{The uses relation is not a data flow diagram.  In the code there will often
be an import statement in module A when it directly uses module B.  Module B
provides the services that module A needs.  The code for module A needs to be
able to see these services (hence the import statement).  Since the uses
relation is transitive, there is a use relation without an import, but the
arrows in the diagram typically correspond to the presence of import statement.}

\wss{If module A uses module B, the arrow is directed from A to B.}

\begin{figure}[H]
\centering
%\includegraphics[width=0.7\textwidth]{UsesHierarchy.png}
\caption{Use hierarchy among modules}
\label{FigUH}
\end{figure}

%\section*{References}

\section{User Interfaces}

\wss{Design of user interface for software and hardware.  Attach an appendix if
needed. Drawings, Sketches, Figma}

\section{Design of Communication Protocols}

\wss{If there is communication between hardware and software, it may be necessary 
to explain the design of the communication protocol.  This section should be filled 
in with NA if it is not appropriate for a given project.}

\section{Database Design}

\wss{If a database is a significant component of the project, the database
design should be summarized here.  This section should be filled in with NA if
it is not appropriate for a given project.}

\section{Timeline}

\wss{Schedule of tasks and who is responsible}

\wss{You can point to GitHub if this information is included there}

\section{Generative AI Use Disclosure}

\wss{ChatGPT (version 5) was used to brainstorm ideas for the template for this disclose section.}

\subsection{AI Tools Used}

\wss{Give the name of the AI tool(s) used and the version.}

\subsection{How and Where Used}

\wss{Explain what the tool(s) were used for.  Examples include: brainstorming,
explaining a technical concept, reviewing the document, generating or modifying
text, generating or modifying diagrams, identifying errors or omissions or
inconsistencies.}

\wss{For each of the uses, identify which parts of the document were affected.}

\wss{You do not need to report every interaction, but you should disclose the 
uses that materially contributed to your work}

\subsection{Verification of AI-Generated Content}

\wss{\begin{itemize}
    \item Describe how the team checked AI-generated or AI-assisted material
    before including it in the document.
    \item In particular, verify factual claims, technical assertions,
    requirements, calculations, references, and other information for which an
    AI system may produce incorrect or fabricated results.
    \item \textbf{The team is responsible for the accuracy, completeness,
consistency, and appropriateness of the submitted document, regardless of
whether material was generated by a human or an AI system.}
\end{itemize}}

\newpage{}

\bibliographystyle {plainnat}
\bibliography{../../../refs/References}

\end{document}